\section{Discussion}
\label{sec:discussion}
\subsection{The security lesson: observable entropy}
The failure is not that the implementation forgets to sample a permutation. It samples one for each block. The failure is that, for a reachable input symmetry, many sampled permutations become observationally equivalent. The security-relevant quantity is the orbit of the transformed object, not the cardinality of the randomizer's support. A $b!$ implementation choice can therefore contribute only $b$ distinguishable calibration states.

This distinction is broader than \scheme{}, but the attack is not a theorem about all linear obfuscation. A designer must ask which transformations survive after quotienting by input stabilizers, which fixed structure remains across observations, and whether recovery metadata becomes an oracle. In the tested construction, the single-row marker is both necessary for legitimate inversion and the label that exposes the columns of $S$. Different marker or mixing structures can change the orbit and the identification problem.

\subsection{Mitigation implications}
Three principles follow. First, do not reuse a transform merely because another component is refreshed: the refreshed component may vanish under a chosen input. Second, bind secrets to the smallest practical identity---request, tenant, or cache object---and specify how shared prefixes and multi-turn sessions inherit that identity. Third, evaluate functional recovery and security together. Replacing a sparse marker with a denser object can defeat this decoder while also breaking the intended inverse; failure of the current attack is not evidence of a correct defense.

Request-specific $S,a$ is the simplest intervention suggested by our measurements. It changes the attacker's problem from amortizing one calibration over other requests to calibrating within the victim's secret scope. Yet the correct engineering choice may instead be authenticated encryption when caches are only stored or transported and need not be manipulated while protected. Obfuscation should be justified by an operation that conventional cryptography cannot efficiently support, not by a general claim that transformed vectors are unintelligible.

\subsection{Limitations and external validity}
Our attacker begins after protected-cache observation is available. We do not show a remote API path to another tenant's cache, measure the prevalence of shared secret epochs, or bypass access control. The study concerns the first Value layer and the tested public implementation. Deeper states, Key recovery, architectures with position-dependent Value inputs, arbitrary orthogonal $M$, and marker patterns outside the analyzed orbit require separate arguments.

The output is unordered within a block. A multiset can expose rare values or keywords, but it is weaker than a verbatim prompt. The synthetic value study uses one model, one secret configuration, 30 cases, and an artificial candidate set. Its incomplete retained traces rule out post-hoc analyses that would require all scores or recovered blocks. The official-entry and multi-epoch studies reuse public conversations; neither supports a claim about an unseen deployment distribution.

The main 99.668\% total weights models with more KV heads and texts with more complete blocks more heavily. It is an audit total, not an estimate over a defined population of users. The model and checkpoint breadth nevertheless tests whether the mechanism is tied to one tokenizer or hidden dimension. Claims about prevalence, per-user harm probability, and production cost remain outside the evidence.

\subsection{Ethics, disclosure, and release}
All experiments ran locally on public model checkpoints, synthetic text, or a publicly released conversation dataset~\cite{lmsys}. We did not recruit or interact with human participants, query a live third-party service, access private caches, or recover non-public user data. No separate human-subjects protocol was submitted because the study involved no participant interaction or new private-data collection. The public conversation subset is used only to exercise file-format and recovery behavior; the manuscript does not reproduce conversation text or identifiers.

The attack is dual use because cache-read access could be converted into content leakage. The authors of \scheme{} were notified of this vulnerability via a GitHub issue on September 28, 2026. The public notice requested a private reporting channel and did not disclose the technical details. We are awaiting a response and will coordinate the detailed disclosure and release timeline with the maintainers. Until affected parties have had a reasonable opportunity to respond, a public artifact should prioritize the proof, frozen aggregate evidence, and inert synthetic reproduction over turnkey integration with live services. No external message is sent by the artifact build process.

We will release an anonymized artifact containing the source snapshot, hashes of frozen result records, figure/table generation, and local synthetic examples, subject to the disclosure timeline. The artifact will distinguish exploratory files from the records supporting paper claims. Any change to a frozen result or denominator will require regenerating the evidence manifest and manuscript assets.
